\documentclass{article}
\usepackage[T1]{fontenc}
\usepackage{lmodern}
\usepackage{graphicx}
\usepackage{amsmath,amssymb}
\usepackage[a4paper,margin=2cm]{geometry}
\usepackage[absolute,overlay]{textpos}
\setlength{\TPHorizModule}{1mm}
\setlength{\TPVertModule}{1mm}
\usepackage[indonesian]{babel}
\usepackage{hyperref}
\setlength{\emergencystretch}{2em}
% unit: Halaman 33

\begin{document}

% === Halaman 33 (llm) ===

\section*{Twofish}

\begin{itemize}
    \item Salah satu finalis AES (\textit{Advanced Encryption Standard})
    \item Dirancang oleh Bruce Schneier.
    \item Twofish dapat dianggap sebagai kelanjutan dari Blowfish
    \item Ukuran blok plainteks dan cipherteks di dalam Twofish adalah 128 bit
    \item Panjang kunci sampai 256 bit
    \item Jumlah putaran 16 kali, setiap putaran melakukan operasi substitusi-permutasi
    \item Jumlah kotak S-box adalah 5 buah
\end{itemize}

\end{document}
